\bibitem[Milovanov(2019)]{ff-milovanov}
Milovanov, A. (2019). \#P-completeness of counting roots of a sparse polynomial. \emph{Information Processing Letters, 142}, 77--79. \url{https://doi.org/10.1016/j.ipl.2018.09.008}. Theorem~1 and Section~3.

\bibitem[Cao et~al.(2026)]{ff-cao}
Cao, Y., Feng, R., Gao, X.-S., Huang, Q.-L., \& Qiu, R. (2026). \emph{Sparse polynomial GCD cannot be computed in fixed-output polynomial time} (Version~2). arXiv:2606.12144v2. \url{https://arxiv.org/abs/2606.12144v2}. Lemma~7, Proposition~9, and equations~(7)--(8), pp.~4--7.

\bibitem[Pries and Zhu(2012)]{ff-pries-zhu}
Pries, R., \& Zhu, H. J. (2012). The $p$-rank stratification of Artin--Schreier curves. \emph{Annales de l'Institut Fourier, 62}(2), 707--726. \url{https://doi.org/10.5802/aif.2692}. Lemma~2.6, p.~712.

\bibitem[A short proof(2026)]{ff-sign}
\emph{A short proof of Deligne's bound for Frobenius eigenvalues, after Alp\"oge}. (2026). Hexagon preprint 2611.00789v1. \url{https://hexagonmath.org/2611.00789v1}. Theorem~1(b), p.~1; Proposition~4, p.~5; Section~8, p.~8. Proof developed by GPT-6 Astra from Alp\"oge's sketch; text by Claude Fable~5.2 using an earlier version by Claude Opus~5.5.

\bibitem[Asok et~al.(2023)]{ff-asok}
Asok, A., Dubouloz, A., \& \O stv\ae r, P. A. (2023). Geometric models for algebraic suspensions. \emph{International Mathematics Research Notices, 2023}(20), 17788--17821. \url{https://doi.org/10.1093/imrn/rnad094}. Theorem~3.1.1, Remark~3.1.3, and Construction~4.2.1; preprint arXiv:2112.08241v2, pp.~8--9 and Section~4.2.

\bibitem[Artin et~al.(1972--1973)]{ff-sga4}
Artin, M., Grothendieck, A., \& Verdier, J.-L. (Eds.). (1972--1973). \emph{Th\'eorie des topos et cohomologie \'etale des sch\'emas (SGA~4)} (Lecture Notes in Mathematics 269, 270, 305). Springer. Expos\'e~XVII, Theorems~5.2.6, 5.4.3; Expos\'e~XVIII, equation~(2.8.1), Theorem~2.9. \url{https://grothendiecksga.com/read/sga4/fr/XVII-5.html}.

\bibitem[Mazza et~al.(2006)]{ff-mvw}
Mazza, C., Voevodsky, V., \& Weibel, C. (2006). \emph{Lecture notes on motivic cohomology} (Clay Mathematics Monographs~2). American Mathematical Society. Example~13.17, Theorem~15.2, Corollary~15.5, and Theorem~19.1. \url{https://sites.math.rutgers.edu/~weibel/MVWnotes/prova-hyperlink.pdf}.

\bibitem[Riou(2010)]{ff-riou}
Riou, J. (2010). Algebraic $K$-theory, $\mathbb A^1$-homotopy and Riemann--Roch theorems. \emph{Journal of Topology, 3}(2), 229--264. \url{https://doi.org/10.1112/jtopol/jtq005}. Theorem~2.3.1, Proposition~2.3.2; preprint arXiv:0907.2710v2, pp.~10--11.

\bibitem[von zur Gathen et~al.(1996)]{ff-gathen}
von zur Gathen, J., Karpinski, M., \& Shparlinski, I. (1996). Counting curves and their projections. \emph{Computational Complexity, 6}, 64--99. Theorem~4.7, Section~4; author's version p.~15. \url{https://cosec.bit.uni-bonn.de/fileadmin/user_upload/publications/pubs/gatkar96.pdf}.

\bibitem[Harvey(2015)]{ff-harvey}
Harvey, D. (2015). Computing zeta functions of arithmetic schemes. \emph{Proceedings of the London Mathematical Society, 111}(6), 1379--1401. \url{https://doi.org/10.1112/plms/pdv056}. Theorem~1.1.

\bibitem[Lauder and Wan(2008)]{ff-lauder-wan}
Lauder, A. G. B., \& Wan, D. (2008). Counting points on varieties over finite fields of small characteristic. In J. P. Buhler \& P. Stevenhagen (Eds.), \emph{Algorithmic number theory: Lattices, number fields, curves and cryptography} (MSRI Publications~44, pp.~579--612). Cambridge University Press. Theorem~1; preprint p.~2. \url{https://arxiv.org/abs/math/0612147}.

\bibitem[Lee and Vakil(2013)]{ff-lee-vakil}
Lee, S. H., \& Vakil, R. (2013). Mn\"ev--Sturmfels universality for schemes. In \emph{A celebration of algebraic geometry} (Clay Mathematics Proceedings~18, pp.~457--468). American Mathematical Society. \url{https://arxiv.org/abs/1202.3934v2}. Theorem~1.3.

\bibitem[Grochow and Qiao(2023)]{ff-grochow-qiao}
Grochow, J. A., \& Qiao, Y. (2023). On the complexity of isomorphism problems for tensors, groups, and polynomials I: Tensor isomorphism-completeness. \emph{SIAM Journal on Computing, 52}(2), 568--617. \url{https://doi.org/10.1137/21M1441110}. Theorem~A, Corollary~A, and Theorem~B, Sections~1 and~3.

\bibitem[Deligne(1974)]{ff-deligne}
Deligne, P. (1974). La conjecture de Weil: I. \emph{Publications Math\'ematiques de l'IH\'ES, 43}, 273--307. \url{https://doi.org/10.1007/BF02684373}. Theorem~1.6, p.~276.

\bibitem[Huybrechts(2023)]{ff-huybrechts}
Huybrechts, D. (2023). \emph{The geometry of cubic hypersurfaces} (Cambridge Studies in Advanced Mathematics~206). Cambridge University Press. \url{https://www.math.uni-bonn.de/people/huybrech/Notes.pdf}. Chapter~1, Corollary~1.12, p.~18, and Section~1.6, p.~27, in the author's edition.

\bibitem[Katz and Rodriguez Villegas(2026)]{ff-katz-rv}
Katz, N. M., \& Rodriguez Villegas, F. (2026). \emph{Remarks on polynomial count varieties}. arXiv:2603.07062v1. \url{https://arxiv.org/abs/2603.07062v1}. Example~2.3, pp.~4--5, and Section~3, pp.~5--6.

\bibitem[Katz(1973)]{ff-sga7}
Katz, N. M. (1973). Une formule de congruence pour la fonction $\zeta$. In P. Deligne \& N. M. Katz (Eds.), \emph{Groupes de monodromie en g\'eom\'etrie alg\'ebrique II (SGA~7~II)} (Lecture Notes in Mathematics~340, Expos\'e~XXII, pp.~401--438). Springer. Theorem~3.1 and Corollary~3.2, p.~410. \url{https://web.math.princeton.edu/~nmk/old/congzeta.pdf}.
